% Frontespizio e retro, sul modello delle tesi seguite da Vardanega.
% Pandoc chiama \maketitle quando metadata.yaml ha un titolo.
\usepackage{graphicx}
\makeatletter
\renewcommand{\maketitle}{%
  \begin{titlepage}
    \centering
    {\Large\bfseries Università di Padova\par}
    \vspace{0.8em}
    {\large\textsc{Dipartimento di Matematica “Tullio Levi-Civita”}\par}
    \vspace{0.8em}
    {\large\textsc{Corso di Laurea in Informatica}\par}
    \vfill
    \includegraphics[width=0.38\textwidth]{figure/unipd-logo.png}\par
    \vfill
    {\LARGE\bfseries \@title\par}
    \vspace{1.2em}
    {\large\scshape Tesi di Laurea\par}
    \vfill
    \begin{minipage}[t]{0.47\textwidth}
      \raggedright
      {\scshape Laureando}\\[0.3em]
      \@author\\[0.3em]
      {\scshape Matricola} 2101055
    \end{minipage}\hfill
    \begin{minipage}[t]{0.47\textwidth}
      \raggedleft
      {\scshape Relatore}\\[0.3em]
      Prof.\ Tullio Vardanega
    \end{minipage}
    \vfill
    \rule{0.8\textwidth}{0.4pt}\par
    \vspace{0.6em}
    {\large\textsc{Anno Accademico 2025-2026}\par}
  \end{titlepage}
  \thispagestyle{empty}
  \vspace*{\fill}
  \noindent\small © \@author. Tutti i diritti riservati.
  Tesi di Laurea: “\@title”, Università di Padova,
  Dipartimento di Matematica “Tullio Levi-Civita”, Corso di Laurea in Informatica.
  \normalsize
  \clearpage
}
\makeatother
